\chapter*{Introducción}
\addcontentsline{toc}{chapter}{Introducción}

% La introducción no cuenta como primer capítulo
\section{Motivación y contexto}
\noindent Cada día ocurren, en promedio, más de 230 accidentes viales en la Ciudad de México, y al menos una persona muere en el lugar del incidente. Estas cifras, que corresponden al promedio entre 2014 y 2023 de los registros del Centro de Comando, Control, Cómputo, Comunicaciones y Contacto Ciudadano (C5), muestran que la siniestralidad vial es un problema de salud pública de primer orden, tanto para la seguridad vial como para el sistema de salud.\\

\indent Más allá del nivel absoluto de incidentes, la evolución reciente de estos eventos sugiere que la situación se ha vuelto especialmente preocupante en los últimos años. La Figura~\ref{fig:incidentes-anuales} muestra la trayectoria del número total de accidentes viales entre 2014 y 2023, desagregados en accidentes menores, con lesionados y fatales. El sistema de fotocívicas entró en vigor el 22 de abril de 2019. Hasta ese momento se observaba una tendencia a la baja tanto en el número total de incidentes como en los accidentes fatales anuales. Sin embargo, tras una caída pronunciada en 2020, muy probablemente asociada a la reducción de la movilidad durante la pandemia de COVID‑19, se registra un repunte en el número de incidentes viales, especialmente en los de mayor severidad.

\begin{figure}
    \centering
    \caption{Número anual de incidentes viales en la CDMX}
    \includegraphics[width=1\linewidth]{Imagenes/yearly-number-of-incidents.pdf}
    \label{fig:incidentes-anuales}
\end{figure}

Este cambio de tendencia se aprecia con claridad al comparar 2018, último año completo previo a la implementación de las fotocívicas, con 2023. En ese periodo, los incidentes totales aumentaron casi 10\%, y los accidentes con al menos una persona fallecida, 57\%. La Figura~\ref{fig:fatales-anuales}, que presenta únicamente el número de accidentes fatales por año, muestra un mínimo en 2018 seguido de un crecimiento sostenido. Aunque con esta evidencia descriptiva no es posible atribuir causalmente estos cambios a un solo factor, sí resulta claro que la introducción del programa de fotocívicas no evitó el repunte observado en los accidentes más graves en el total agregado de la ciudad.

\begin{figure}
    \centering
    \caption{Número anual de accidentes fatales}
    \includegraphics[width=1\linewidth]{Imagenes/yearly-fatal-number-of-incidents.pdf}
    \label{fig:fatales-anuales}
\end{figure}

La composición de los siniestros también se ha modificado. Como se observa en la Figura~\ref{fig:incidentes-anuales}, la disminución de 2020 en el número total de incidentes se explica principalmente por la reducción de los accidentes menores, mientras que los accidentes con lesionados o fallecidos no se reducen en la misma proporción. A partir de 2020, los incidentes con personas lesionadas superan a los accidentes considerados menores y se mantienen por encima hasta 2023, ampliando cada año la brecha entre ambos tipos de eventos. Es decir, no solo sigue habiendo muchos accidentes, sino que una fracción creciente de ellos tiene consecuencias más severas.\\

Este diagnóstico coincide con lo reportado por la Secretaría de Movilidad (SEMOVI) en el Programa Integral de Seguridad Vial de la Ciudad de México 2021–2024 (SEMOVI, 2021), donde se documenta que la reducción en la afluencia vehicular durante la pandemia, de aproximadamente 45\%, permitió mayores velocidades de circulación y que el aumento relativo de incidentes con lesionados y fallecidos se concentró en vías primarias y de acceso controlado. En otras palabras, el propio gobierno local reconoce que la velocidad y las características de la infraestructura vial juegan un papel central en la siniestralidad reciente.\\ 

\section{Programas de control de velocidad en la Ciudad de México}
\noindent El control de la velocidad ha adquirido relevancia mundial como estrategia central para mejorar la seguridad vial. La Guía para el Control de la Velocidad (García Cruz, 2018) identifica la velocidad como el principal factor de riesgo en la ocurrencia y severidad de los siniestros viales y destaca que “el 90\% de los automovilistas se considera un conductor por encima del promedio y de bajo riesgo”. Esta sobreconfianza convive con una percepción social de la velocidad como algo valioso, también en el transporte, porque se asocia con ahorrar tiempo, aunque en zonas urbanas ese ahorro no sea significativo. La misma guía documenta que la probabilidad de sobrevivir a un atropellamiento se reduce de 90\% a 30 km/h a poco más de 10\% a 50 km/h, lo que refuerza la necesidad de limitar las velocidades efectivas en la red vial.\\

\indent En la Ciudad de México, la principal herramienta de control automatizado de la velocidad han sido los programas de fotomultas y, posteriormente, de fotocívicas. El Programa Fotomultas comenzó a operar en diciembre de 2015 y sancionaba el exceso de velocidad detectado por los radares con multas económicas para el propietario del vehículo. Al sustituirlo, el propio gobierno de la ciudad lo describió como una política con fines recaudatorios e ineficaz para cambiar conductas, en la que, además, la ubicación de las cámaras no era pública (SEMOVI, 2019).\\

En abril de 2019, el esquema de fotomultas se sustituyó por el sistema de fotocívicas, cuyo objetivo declarado es cambiar la conducta de los conductores y no recaudar. El nuevo modelo se basa en un sistema de puntos asociado a las placas de los vehículos registrados en la ciudad. Cada placa inicia con diez puntos y, por cada infracción detectada, se pierde un punto (cinco puntos si el límite de velocidad se excede en más de 40\%). La forma de resarcir estas infracciones ya no es el pago de una multa, sino la realización de cursos educativos y trabajo comunitario. Cuando una placa pierde más de dos puntos, su propietario no puede verificar el vehículo hasta cumplir con las sanciones correspondientes (no es posible circular si el vehículo no está verificado). Para los vehículos con placas foráneas, en cambio, el esquema mantiene las sanciones económicas.\\

La implementación de fotocívicas vino acompañada de una reubicación de los radares. De acuerdo con la SEMOVI (2019), las cámaras se reubicaron hacia tramos con mayor incidencia de hechos de tránsito, en particular aquellos con más muertos y lesionados, así como en puntos con elevados volúmenes de tránsito y donde el límite de velocidad tendía a no respetarse. En total se instalaron 113 radares. Aunque la lógica oficial no menciona explícitamente indicadores como el exceso medio de velocidad, sí enfatiza la concentración de incidentes y la importancia de las vías primarias y de acceso controlado como criterio de ubicación.\\

En la práctica, sin embargo, el nuevo sistema ha seguido funcionando en gran medida como un esquema de sanciones económicas. Los registros de la propia autoridad (SSC, 2025), que se presentan en la Figura~\ref{fig:sanciones}, muestran que la proporción de sanciones cívicas ha oscilado entre 14\% y 24\% del total, lo que implica que la mayoría de las infracciones se ha traducido en multas monetarias. Esto puede deberse, entre otras razones, al elevado número de vehículos con placas foráneas que circulan por la ciudad o a que los propietarios con placas locales internalizan mejor el riesgo de sanción y ajustan su comportamiento. Aunque el análisis detallado de estas causas excede el alcance de esta tesis, el dato es relevante porque sugiere que la experiencia de los conductores con el sistema se asemeja más a un esquema de multas que al modelo cívico originalmente planteado.\\

\begin{figure}
    \centering
    \caption{Total de infracciones por tipo de sanción}
    \includegraphics[width=1\linewidth]{Imagenes/yearly-number-of-sanctions.pdf}
    \label{fig:sanciones}
\end{figure}

\section{Pregunta de investigación e hipótesis}
\noindent En este contexto, es natural preguntarse si las cámaras de velocidad realmente han tenido el efecto que prometían sobre la seguridad vial de la ciudad. A primera vista, podría pensarse que, al menos en los puntos donde se instalaron, deberían observarse menos incidentes, en particular menos hechos con personas lesionadas o fallecidas. Sin embargo, las tendencias mostradas en las figuras anteriores, así como el repunte reciente de los accidentes fatales, hacen que esta intuición deje de ser tan obvia. La pregunta que guía esta tesis es, precisamente, \textbf{¿cómo han cambiado los incidentes viales alrededor de las cámaras de velocidad desde la implementación del sistema de fotocívicas?}\\

La hipótesis con la que se empezó este análisis es que las cámaras reducen los incidentes viales en su entorno, con un efecto mayor en los incidentes con lesionados o fatalidades, y en magnitudes similares a las que reporta la literatura internacional, de entre 10\% y 20\%.

\section{Diseño y resultados}
\noindent Para responder la pregunta, se construyó un diseño que busca comparar zonas que realmente sean comparables. Alrededor de las cámaras se trazaron círculos de entre 100 y 300 metros, y cada uno se emparejó por \emph{propensity score} con un círculo del mismo tamaño, sobre la misma red de vías primarias, que antes del programa era igual de peligroso. Con los incidentes registrados por el C5, se estimó un modelo de diferencias en diferencias que compara cómo cambiaron los incidentes en ambos grupos después de abril de 2019. Como la ventana de análisis incluye la pandemia, se revisó también, año por año, que esta no cargara el resultado.\\

Las cámaras no muestran un efecto decisivo sobre el número de incidentes viales. Después del programa, los incidentes bajaron casi lo mismo en las zonas con cámara que en sus controles, y en ningún radio ni tipo de incidente la diferencia se distingue de cero. Lo que sí se puede decir es que, si las cámaras redujeron los incidentes, lo hicieron mucho menos de lo esperado, ya que los datos descartan las reducciones de 10\% a 20\% que planteaba la hipótesis. En los incidentes fatales, en cambio, no se puede concluir nada, porque son tan pocos que el efecto tendría que haber sido muy grande para distinguirse del ruido.\\

El resultado depende, sobre todo, de contra qué se comparen las cámaras. Si las zonas con cámara se comparan consigo mismas antes y después del programa, o contra cualquier zona sin cámara de la misma red vial, los incidentes caen alrededor de 14\%, una reducción similar a la de la literatura. Esa reducción desaparece al compararlas contra zonas que antes del programa eran igual de peligrosas. Como las cámaras se ubicaron donde ya había más incidentes, las comparaciones más sencillas les atribuyen una caída que también vivieron las zonas parecidas que no tenían cámara.

\section{Contribuciones}
\noindent Esta tesis se construye sobre los dos trabajos que ya han evaluado las cámaras de velocidad en la Ciudad de México. Quintero et al.\ (2023) estudian el efecto del programa a nivel alcaldía, y Madrigal y Rodríguez (2021), en el tramo inmediato a cada cámara, con cruces de la red primaria sin emparejar como controles. Frente a ellos, esta tesis aporta tres cosas.\\

La primera es un contrafactual construido para cada zona con cámara. Como las cámaras no se ubicaron al azar, cada zona con cámara se compara con otra de la misma red vial que antes del programa tenía una siniestralidad, una infraestructura vial y una afluencia del Metro parecidas, y no con cualquier vialidad sin cámara. Además, como los pares se construyen con los mismos datos con los que se estima el efecto, la incertidumbre se calcula repitiendo todo el proceso, desde el \emph{propensity score}.\\

La segunda es que el efecto se estima en cinco radios con las mismas unidades tratadas. Así, las diferencias entre un radio y otro se deben al tamaño del área que se analiza, y no a que cambie el conjunto de cámaras que se compara.\\

La tercera es mostrar, con los mismos datos, cuánto cambia la conclusión según la comparación que se haga, lo que ayuda a entender por qué distintas evaluaciones de un mismo programa pueden llegar a conclusiones distintas.\\

El resto de la tesis se organiza de la siguiente manera: el capítulo 1 revisa la evidencia internacional sobre cámaras de velocidad y los estudios previos para la Ciudad de México; el capítulo 2 presenta la estrategia empírica y los supuestos bajo los cuales los resultados pueden interpretarse como causales; el capítulo~\ref{chap:datos} describe los datos, la construcción de las unidades de análisis y el emparejamiento; el capítulo~\ref{chap:resultados} presenta los resultados, y, por último, las conclusiones discuten qué se puede decir a partir de ellos, sus limitaciones y las posibles explicaciones de lo que se encontró.
